\documentclass{beamer}

\usepackage[T1]{fontenc}
\usepackage[utf8]{inputenc}
\usepackage{lmodern}
\usepackage{amsmath,amssymb,amsfonts}
\usepackage{physics}
\usepackage{bm}
\usepackage{mathtools}
\usepackage{quantikz}
\usepackage{listings}
\usepackage{xcolor}

\lstset{
  basicstyle=\ttfamily\tiny,
  keywordstyle=\color{blue},
  commentstyle=\color{green!50!black},
  showstringspaces=false,
  breaklines=true
}

\title{Using QAOA to Optimize Mean Squared Error}
\subtitle{From MSE to QUBO, Ising Hamiltonians, and Variational Quantum Circuits}
\author{}
\date{}

\begin{document}

\frame{\titlepage}

\begin{frame}{Outline}
\begin{enumerate}
\item Why MSE is not directly a QAOA problem
\item Binary encoding of continuous parameters
\item MSE \(\rightarrow\) QUBO \(\rightarrow\) Ising mapping
\item Fully explicit one-parameter linear fit
\item Extension to two-parameter regression
\item Regularized regression: \(L_2\) and \(L_1\)-type penalties
\item PDE discretizations and QAOA
\item QAOA circuit and Qiskit implementation sketch
\item Limitations and outlook
\end{enumerate}
\end{frame}

\begin{frame}{Why QAOA Is Not Directly an MSE Optimizer}
The mean squared error is usually a continuous objective:
\[
L(\theta)=\frac1N \sum_{i=1}^N \left(y_i-f(x_i;\theta)\right)^2.
\]

QAOA is designed for discrete optimization:
\[
\min_{z\in\{0,1\}^m} C(z),
\]
or equivalently for a cost Hamiltonian diagonal in the computational basis.

Therefore the strategy is:
\[
\text{continuous MSE}
\;\longrightarrow\;
\text{binary encoding}
\;\longrightarrow\;
\text{QUBO}
\;\longrightarrow\;
\text{Ising Hamiltonian}
\;\longrightarrow\;
\text{QAOA}.
\]
\end{frame}

\begin{frame}{Binary Encoding of Parameters}
Suppose a parameter \(\theta\) is restricted to a finite interval and represented by \(m\) bits:
\[
\theta = \theta_{\min} + \Delta \sum_{k=0}^{m-1} 2^k z_k,
\qquad z_k\in\{0,1\}.
\]

Here:
\begin{itemize}
\item \(\theta_{\min}\): lower bound,
\item \(\Delta\): resolution,
\item \(z_k\): binary decision variables.
\end{itemize}

For several parameters \(\theta_1,\dots,\theta_p\), we encode each parameter separately.
Once substituted into the MSE, the loss becomes a polynomial in binary variables.
\end{frame}

\begin{frame}{From MSE to QUBO}
A quadratic unconstrained binary optimization (QUBO) problem has the form
\[
C(z)=c_0+\sum_i c_i z_i+\sum_{i<j} c_{ij} z_i z_j,
\qquad z_i\in\{0,1\}.
\]

If the model is linear in the encoded binary variables, then the squared loss naturally becomes quadratic in the bits.

This is ideal for QAOA, because QUBO maps directly to an Ising Hamiltonian.
\end{frame}

\begin{frame}{QUBO to Ising Mapping}
Use the standard relation
\[
z_i=\frac{1-Z_i}{2},
\]
where \(Z_i\) is the Pauli-\(Z\) operator on qubit \(i\).

Then a QUBO
\[
C(z)=c_0+\sum_i c_i z_i+\sum_{i<j} c_{ij} z_i z_j
\]
becomes
\[
H_C = \text{const} + \sum_i h_i Z_i + \sum_{i<j} J_{ij} Z_i Z_j.
\]

This Hamiltonian is diagonal in the computational basis and is therefore a valid QAOA cost Hamiltonian.
\end{frame}

\begin{frame}{Worked Example: One-Parameter Linear Fit}
Consider the model
\[
f(x;a)=ax,
\]
and the dataset
\[
(x_1,y_1)=(1,1), \qquad (x_2,y_2)=(2,2).
\]

Take a two-bit encoding
\[
a=z_0+2z_1,
\qquad z_0,z_1\in\{0,1\}.
\]

Then
\[
L(a)=\frac12\left[(1-a)^2 + (2-2a)^2\right].
\]

We now express this entirely in terms of \(z_0,z_1\).
\end{frame}

\begin{frame}{Expanding the Worked Example}
Insert
\[
a=z_0+2z_1
\]
into
\[
L(a)=\frac12\left[(1-a)^2+(2-2a)^2\right].
\]

Since
\[
(1-a)^2 = 1 - 2a + a^2,
\qquad
(2-2a)^2 = 4 - 8a + 4a^2,
\]
we get
\[
L(a)=\frac12\left(5 - 10a + 5a^2\right)
=
\frac52 - 5a + \frac52 a^2.
\]

Now substitute \(a=z_0+2z_1\).
\end{frame}

\begin{frame}{Compute the Binary Polynomial}
We have
\[
a^2 = (z_0+2z_1)^2 = z_0^2 + 4z_1^2 + 4z_0z_1.
\]

Because \(z_i^2=z_i\) for binary variables,
\[
a^2 = z_0 + 4z_1 + 4z_0z_1.
\]

Therefore
\[
L(z_0,z_1)
=
\frac52 - 5(z_0+2z_1) + \frac52(z_0+4z_1+4z_0z_1).
\]

Collecting terms:
\[
\boxed{
L(z_0,z_1)=\frac52 - \frac52 z_0 + 10 z_0z_1
}
\]

Remarkably, the \(z_1\) term cancels in this particular dataset.
\end{frame}

\begin{frame}{QUBO Coefficients for the Worked Example}
Thus the QUBO is
\[
C(z_0,z_1)=c_0+c_1 z_0 + c_2 z_1 + c_{12} z_0 z_1
\]
with
\[
c_0=\frac52,\qquad
c_1=-\frac52,\qquad
c_2=0,\qquad
c_{12}=10.
\]

This is now in a form directly usable for QAOA.
\end{frame}

\begin{frame}{Map the Worked Example to an Ising Hamiltonian}
Use
\[
z_0=\frac{1-Z_0}{2},\qquad
z_1=\frac{1-Z_1}{2}.
\]

Then
\[
z_0 z_1=\frac14(1-Z_0-Z_1+Z_0Z_1).
\]

So
\[
H_C
=
\frac52
-\frac52 \frac{1-Z_0}{2}
+10\cdot \frac14(1-Z_0-Z_1+Z_0Z_1).
\]

Simplify term by term.
\end{frame}

\begin{frame}{Final Cost Hamiltonian}
After simplification:
\[
H_C
=
\frac{15}{4}
-\frac{5}{4} Z_0
-\frac{5}{2} Z_1
+\frac{5}{2} Z_0 Z_1.
\]

Thus the corresponding Ising coefficients are
\[
\text{const}=\frac{15}{4},\qquad
h_0=-\frac54,\qquad
h_1=-\frac52,\qquad
J_{01}=\frac52.
\]

This is the Hamiltonian used in the QAOA cost unitary
\[
U_C(\gamma)=e^{-i\gamma H_C}.
\]
\end{frame}

\begin{frame}{Check the Four Bitstrings}
The bitstrings correspond to:
\[
00 \Rightarrow a=0,\qquad
01 \Rightarrow a=2,\qquad
10 \Rightarrow a=1,\qquad
11 \Rightarrow a=3.
\]

Evaluate the MSE:
\[
L(0)=\frac12(1^2+2^2)=\frac52,
\]
\[
L(1)=0,
\]
\[
L(2)=\frac12(1+4)=\frac52,
\]
\[
L(3)=\frac12(4+16)=10.
\]

So the optimum is clearly at
\[
a=1 \quad \Longleftrightarrow \quad z_0=1,\; z_1=0.
\]
\end{frame}

\begin{frame}{QAOA Ansatz for the Worked Example}
For two qubits and depth \(p=1\), the QAOA state is
\[
\ket{\psi(\gamma,\beta)}
=
e^{-i\beta(X_0+X_1)} e^{-i\gamma H_C} \ket{+}^{\otimes 2}.
\]

The standard mixer Hamiltonian is
\[
H_M = X_0 + X_1.
\]

The optimization target is
\[
\min_{\gamma,\beta} \bra{\psi(\gamma,\beta)} H_C \ket{\psi(\gamma,\beta)}.
\]
\end{frame}

\begin{frame}{QAOA Circuit for the Worked Example}
\centering
\begin{quantikz}
\lstick{$\ket{+}$} & \gate{$e^{-i\gamma H_C}$} & \gate{$R_x(2\beta)$} & \meter{} \\
\lstick{$\ket{+}$} & \gate{$e^{-i\gamma H_C}$} & \gate{$R_x(2\beta)$} & \meter{}
\end{quantikz}

\vspace{0.7em}

The cost unitary itself can be decomposed into
\begin{itemize}
\item one-qubit \(R_z\) rotations from the \(Z_i\) terms,
\item a two-qubit entangling gate from the \(Z_0 Z_1\) term.
\end{itemize}
\end{frame}

\begin{frame}{Two-Parameter Regression}
Now consider
\[
f(x;a,b)=ax+b.
\]

The MSE is
\[
L(a,b)=\frac1N\sum_{i=1}^N \left(y_i-a x_i-b\right)^2.
\]

Encode
\[
a=a_{\min}+\sum_{k=0}^{m_a-1}\alpha_k z_k,
\qquad
b=b_{\min}+\sum_{\ell=0}^{m_b-1}\beta_\ell w_\ell.
\]

Then the loss becomes a quadratic polynomial in all binary variables
\[
q=(z_0,\dots,z_{m_a-1},w_0,\dots,w_{m_b-1}).
\]
\end{frame}

\begin{frame}{General Two-Parameter QUBO Structure}
Since the model is affine in \(a\) and \(b\), the squared loss produces terms of the form
\[
L(q)
=
c_0 + \sum_i c_i q_i + \sum_{i<j} c_{ij} q_i q_j.
\]

The linear coefficients come from:
\begin{itemize}
\item \(a\)-terms,
\item \(b\)-terms,
\item offsets from \(a_{\min},b_{\min}\).
\end{itemize}

The quadratic coefficients come from:
\begin{itemize}
\item \(a^2\),
\item \(b^2\),
\item \(ab\)-cross terms.
\end{itemize}

This remains a QUBO problem and therefore admits a direct Ising representation.
\end{frame}

\begin{frame}{Regularized Regression: \(L_2\) Penalty}
Suppose we add ridge regularization:
\[
L_{\mathrm{ridge}}(a,b)
=
\frac1N\sum_{i=1}^N (y_i-ax_i-b)^2 + \lambda(a^2+b^2).
\]

After binary encoding, the \(a^2\) and \(b^2\) terms remain quadratic in the bits.

Hence:
\[
\boxed{
\text{MSE + } L_2 \text{ regularization still maps naturally to QUBO.}
}
\]

So ridge-type penalties are straightforward to incorporate into QAOA.
\end{frame}

\begin{frame}{Regularized Regression: \(L_1\)-Type Penalty}
If instead we add
\[
L_{\mathrm{lasso}}(a,b)
=
\frac1N\sum_{i=1}^N (y_i-ax_i-b)^2 + \lambda(|a|+|b|),
\]
then the absolute value is not polynomial.

Possible remedies:
\begin{itemize}
\item discretize \(a,b\) to be nonnegative,
\item introduce auxiliary binary variables,
\item represent sign and magnitude separately,
\item use penalty constructions to approximate \(|a|\).
\end{itemize}

Thus \(L_1\)-type penalties can be handled, but the QUBO construction becomes more elaborate.
\end{frame}

\begin{frame}{From PDEs to Optimization}
Many PDEs become linear systems after discretization:
\[
A\mathbf{u}=\mathbf{f}.
\]

Equivalently, one can solve the variational problem
\[
\min_{\mathbf{u}}
\left(
\frac12 \mathbf{u}^T A \mathbf{u} - \mathbf{f}^T \mathbf{u}
\right).
\]

If \(\mathbf{u}\) is treated as continuous, this is not a natural QAOA problem.

But if the field values are discretized into binary variables, then QAOA can be used on the resulting QUBO.
\end{frame}

\begin{frame}{Binary Encoding of a PDE Field}
Suppose a discretized field is represented by \(u_i\), and each \(u_i\) is encoded as
\[
u_i = u_{\min} + \Delta \sum_{k=0}^{m-1} 2^k z_{ik}.
\]

Then the variational PDE functional becomes a polynomial in the binary variables \(z_{ik}\).

For linear elliptic problems, this typically yields a quadratic form:
\[
C(z)=z^T Q z + c^T z + c_0,
\]
that is, again a QUBO.

So in principle:
\[
\text{discretized PDE} \;\longrightarrow\; \text{binary field encoding} \;\longrightarrow\; \text{QAOA}.
\]
\end{frame}

\begin{frame}{Interpretation for PDEs}
This is conceptually possible, but usually not the most efficient approach.

Why?
\begin{itemize}
\item PDEs are naturally continuous problems,
\item fine discretization requires many bits,
\item the number of qubits grows rapidly,
\item standard linear algebra or VQE-like continuous variational approaches are often better suited.
\end{itemize}

So QAOA is more natural for PDE-inspired problems with
\begin{itemize}
\item discrete controls,
\item combinatorial constraints,
\item or low-resolution binary field models.
\end{itemize}
\end{frame}

\begin{frame}{Qiskit Sketch: Build the Cost Hamiltonian}
Below is a minimal sketch for the worked two-qubit example.

\begin{lstlisting}[language=Python]
from qiskit.quantum_info import SparsePauliOp

# H_C = 15/4 - 5/4 Z0 - 5/2 Z1 + 5/2 Z0 Z1
cost_hamiltonian = SparsePauliOp.from_list([
    ("II", 15/4),
    ("ZI", -5/4),
    ("IZ", -5/2),
    ("ZZ",  5/2),
])
\end{lstlisting}

This Hamiltonian can be fed into a QAOA routine together with a classical optimizer.
\end{frame}

\begin{frame}{Qiskit Sketch: Run QAOA}
\begin{lstlisting}[language=Python]
from qiskit_algorithms import QAOA
from qiskit_algorithms.optimizers import COBYLA
from qiskit.primitives import Sampler

sampler = Sampler()
optimizer = COBYLA(maxiter=100)

qaoa = QAOA(
    sampler=sampler,
    optimizer=optimizer,
    reps=1
)

result = qaoa.compute_minimum_eigenvalue(cost_hamiltonian)
print(result.eigenvalue)
print(result.best_measurement)
\end{lstlisting}

The best measured bitstring is then decoded back into the model parameter \(a\).
\end{frame}

\begin{frame}{Decoding the Bitstring}
For the worked example, the decoding rule is
\[
a=z_0+2z_1.
\]

Thus:
\begin{itemize}
\item bitstring \(10\) means \(a=1\),
\item bitstring \(01\) means \(a=2\),
\item bitstring \(11\) means \(a=3\).
\end{itemize}

After the QAOA optimization, one takes the highest-probability or lowest-energy bitstring and converts it back into the regression parameters.
\end{frame}

\begin{frame}{Advantages and Limitations}
Advantages:
\begin{itemize}
\item natural for discrete or combinatorial regression variants,
\item can incorporate binary constraints directly,
\item unified QUBO framework.
\end{itemize}

Limitations:
\begin{itemize}
\item continuous MSE must be discretized,
\item encoding overhead can be large,
\item standard gradient-based regression is usually much more efficient,
\item current hardware limits practical scale.
\end{itemize}
\end{frame}

\begin{frame}{When QAOA Is Most Natural}
QAOA is most appropriate when the learning problem has discrete structure, for example:
\begin{itemize}
\item sparse feature selection,
\item binary model selection,
\item quantized parameter search,
\item discrete control optimization,
\item low-resolution binary PDE field approximations,
\item regression with combinatorial penalties.
\end{itemize}

For smooth continuous least-squares problems, classical optimization or VQE-like continuous methods are usually more natural.
\end{frame}

\begin{frame}{Summary}
To use QAOA for MSE minimization:
\begin{enumerate}
\item write the MSE,
\item encode parameters with bits,
\item expand the result as a QUBO,
\item map the QUBO to an Ising Hamiltonian,
\item run QAOA,
\item decode the optimal bitstring into the fitted parameters.
\end{enumerate}

In the fully explicit straight-line example,
\[
L(z_0,z_1)=\frac52 - \frac52 z_0 + 10 z_0 z_1
\]
and the QAOA cost Hamiltonian is
\[
H_C
=
\frac{15}{4}
-\frac{5}{4} Z_0
-\frac{5}{2} Z_1
+\frac{5}{2} Z_0 Z_1.
\]
\end{frame}

\begin{frame}{Possible Extensions}
Natural follow-up topics include:
\begin{itemize}
\item explicit \(a,b\) regression with four or more qubits,
\item higher-order polynomial fitting,
\item constrained regression and feature selection,
\item PDE discretizations leading to QUBOs,
\item comparison with VQE and quantum linear-system solvers.
\end{itemize}
\end{frame}

\end{document}
